

El presente capítulo establece los criterios, procedimientos y responsabilidades que rigen la determinación de aeronavegabilidad de los \textbf{Sistemas de Aeronaves No Tripuladas (UAS)} de SKYFILMS, asegurando que toda aeronave sea operada únicamente cuando cumpla las condiciones técnicas, funcionales y documentales exigidas por los Reglamentos Aeronáuticos de Colombia (RAC 100) y por las disposiciones de la Aerocivil.\\

El objetivo es garantizar que ningún UAS se incorpore a la operación si no ha sido inspeccionado, verificado y certificado como aeronavegable, preservando la seguridad operacional, la trazabilidad del mantenimiento y la integridad del sistema.

\section{Alcance}
Este procedimiento aplica a todas las aeronaves, componentes y sistemas UAS registrados bajo el \textbf{Certificado de Explotador de SKYFILMS}, y cubre:

\begin{itemize}
    \item La evaluación técnica de aeronavegabilidad posterior a mantenimiento, inspección o reparación.

    \item La identificación y clasificación de una aeronave como \textbf{aeronavegable o no aeronavegable.}

    \item El procedimiento de reincorporación al servicio luego de la corrección de discrepancias o fallas.
\end{itemize}

Aplica tanto a \textbf{mantenimientos ejecutados internamente} como a los realizados por \textbf{proveedores o talleres externos autorizados}.

\section{Criterios para la Determinación de Aeronavegabilidad}

La aeronavegabilidad de un UAS se determina mediante una evaluación integral que considera tanto los aspectos técnicos como la documentación de respaldo. Los criterios principales son los siguientes:

\begin{enumerate}
    \item \textbf{Cumplimiento del Programa de Mantenimiento:} La aeronave debe haber completado satisfactoriamente todas las inspecciones y tareas preventivas o correctivas establecidas en el programa técnico de SKYFILMS.

    \item \textbf{Conformidad con el diseño y configuración aprobada:} Se verificará que el UAS mantenga las especificaciones originales del fabricante, sin modificaciones no documentadas ni alteraciones que afecten la seguridad.

    \item \textbf{Ausencia de defectos o discrepancias:} No deben existir daños estructurales, fisuras, corrosión, fallas electrónicas o inconsistencias en los sistemas de propulsión, control o comunicación.

    \item \textbf{Documentación técnica completa y vigente:} Los registros, bitácoras y formularios de mantenimiento deben estar debidamente diligenciados, firmados y actualizados.

    \item \textbf{Resultados satisfactorios de pruebas funcionales:} Todas las pruebas operacionales exigidas por el fabricante y el Técnico UAS deben haber sido ejecutadas y aprobadas sin observaciones.

    \item \textbf{Autorización del personal competente:} Solo el personal técnico designado (Técnico UAS) podrá certificar oficialmente la condición aeronavegable.
\end{enumerate}

Estos criterios aseguran uniformidad y rigor en la determinación de la condición técnica de cada aeronave.

\section{Procedimiento para la Determinación de Aeronavegabilidad}

Este procedimiento define las etapas que deben cumplirse para determinar la condición \textbf{aeronavegable o no aeronavegable} de una aeronave UAS.\\

Cada fase del proceso incluye verificaciones técnicas y documentales que garantizan la trazabilidad y validez del resultado.

\begin{enumerate}
    \item \textbf{Finalización del mantenimiento:} Una vez completadas las tareas técnicas, el Técnico UAS documentará las actividades realizadas, los componentes reemplazados y las observaciones pertinentes en el \textbf{Registro de Mantenimiento}.

    \item \textbf{Inspección técnica post-mantenimiento:} El Técnico de UAS efectuará una revisión visual y funcional de la aeronave, comprobando que el trabajo se haya ejecutado conforme al MCM y al manual del fabricante.

    \item \textbf{Evaluación documental:} Se revisará que toda la información técnica (formatos y reportes de prueba) se encuentre completa, legible y firmada.

    \item \textbf{Pruebas funcionales:} Cuando aplique, se realizarán pruebas en tierra o vuelos de verificación supervisados, verificando la correcta respuesta de los sistemas de control, propulsión y telemetría.

    \item \textbf{Determinación final:}
        \begin{itemize}
            \item Si el UAS cumple todos los requisitos, se declara \textbf{AERONAVEGABLE}.

            \item Si no cumple alguno, se clasifica como \textbf{NO AERONAVEGABLE}, emitiéndose un informe de discrepancias con las acciones correctivas a ejecutar.
        \end{itemize}

    \item \textbf{Certificación:} La aeronavegabilidad se formaliza mediante la emisión del \textbf{Certificado de Aeronavegabilidad Operativa}, firmado por el \textbf{Técnico UAS}.
\end{enumerate}

Este proceso garantiza la transparencia y control técnico en todas las decisiones relacionadas con la aptitud de vuelo de los UAS.

\section{Procedimiento de Reincorporación al Servicio}
Cuando una aeronave haya sido declarada no aeronavegable, deberá seguir un proceso de reincorporación para volver al estado operativo.\\

Este procedimiento permite validar que las correcciones implementadas fueron eficaces y que la aeronave puede retornar al servicio de manera segura:

\begin{enumerate}
    \item \textbf{Ejecución de correcciones:} El Técnico UAS ejecutará las acciones correctivas recomendadas en el informe de discrepancias, registrando los trabajos realizados y los repuestos utilizados.

    \item \textbf{Verificación técnica:} Se inspeccionará nuevamente la aeronave posterior a las correcciones para confirmar que todas las fallas hayan sido solucionadas satisfactoriamente.

    \item \textbf{Pruebas operativas:} Se realizarán las pruebas funcionales o vuelos de comprobación requeridos para validar la operación del sistema.

    \item \textbf{Aprobación de reincorporación:} Una vez superadas las pruebas, el Tecnico UAS emitirá la autorización formal de reincorporación al servicio, levantando la restricción previa.

    \item \textbf{Actualización documental:} El estatus de la aeronave se actualizará en el \textbf{Libro de Mantenimiento}, conservando todos los registros de respaldo.
\end{enumerate}

Este procedimiento garantiza que toda reincorporación se realice bajo control técnico, asegurando la seguridad operacional antes de reanudar vuelos.

\section{Responsabilidades}
La correcta aplicación de este procedimiento depende del compromiso y la ejecución coordinada del personal técnico.\\

Las responsabilidades se definen de la siguiente manera:

\begin{itemize}
    \item \textbf{Gerente General:} Asegura el cumplimiento de las disposiciones técnicas y provee los recursos necesarios para mantener la flota en condiciones aeronavegables.

    \item \textbf{Tecnico UAS:}
        \begin{itemize}
            \item Supervisa el proceso de evaluación técnica, aprueba las determinaciones finales y autoriza la reincorporación de aeronaves al servicio.

            \item Ejecuta inspecciones, verifica registros, realiza pruebas funcionales y conserva la documentación técnica.

            \item Ejecutan las tareas preventivas y correctivas conforme a las especificaciones del fabricante, reportan discrepancias y registran adecuadamente los trabajos realizados.

        \end{itemize}

\end{itemize}

Cada función es esencial para mantener la integridad del proceso y asegurar la validez de las determinaciones ante la autoridad aeronáutica.

\section{Disposiciones Finales}
Toda aeronave operada por SKYFILMS deberá contar con una determinación vigente de aeronavegabilidad antes de su uso operacional. El incumplimiento de este requisito constituirá una no conformidad técnica grave, que podrá conllevar la suspensión temporal de operaciones o la apertura de acciones correctivas inmediatas.\\

La aplicación sistemática de este procedimiento garantiza que todas las aeronaves mantengan su integridad, rendimiento y seguridad conforme a los más altos estándares de la Aerocivil y del Programa de Mantenimiento de SKYFILMS, fortaleciendo la cultura de seguridad operacional dentro de la organización.